\chapter{Implementacja i uruchomienie}
Instrukcja zawiera informacje potrzebne do instalacji i użytkowania aplikacji.

\section{Wymagania sprzętowe}
Po stronie użytkownika:
\begin{itemize}
	\item Smartfon z Android 8.0+
	\item Połączenie Wi-Fi \\
\end{itemize}

Po stronie serwera:
\begin{itemize}
	\item Komputer z XAMPP
	\item ESP32 WROOM
	\item Czujnik DHT11
	\item Zasilanie 5V
\end{itemize}

\section{Instrukcja instalacji}
\begin{itemize}
	\item Połącz komputer (serwer) i telefon do tej samej sieci Wi-Fi.
	\item W pliku MainActivity.java w Android Studio \\ wpisz to IP w linii \texttt{.baseUrl("http://:8080/")} jeśli używasz emulatora nic nie zmieniaj.
	\item Uruchom bazę danych Docker/XAMPP oraz aplikację SpringBoot.
	\item Podłącz telefon do komputera i w Android Studio kliknij zieloną strzałkę Run.
\end{itemize}

\section{Instrukcja użytkowania}
Po uruchomieniu aplikacji użytkownik otrzymuje natychmiastowy dostęp do aktualnej temperatury oraz chronologicznej historii wszystkich zapisów. Dane w programie aktualizują się automatycznie co dziesięć sekund, co zapewnia stały podgląd na najświeższe odczyty bez konieczności ręcznego odświeżania ekranu. System inteligentnie zarządza baterią, wstrzymując pobieranie informacji w momencie zminimalizowania aplikacji i wznawiając pracę zaraz po jej ponownym otwarciu.
